\documentclass[12pt]{article}
\usepackage[T5]{fontenc}
\usepackage[utf8]{inputenc}
\usepackage[vietnamese]{babel}
\shorthandoff{"}

% Cấu hình Font chữ và Cỡ chữ 13pt tiêu chuẩn
\usepackage[fontsize=13pt]{scrextend}
\usepackage{mathptmx}

% Cấu hình lề trang theo quy chuẩn kỹ thuật
\usepackage[top=2.5cm,bottom=2.5cm,left=3.5cm,right=2.0cm]{geometry}

% Cấu hình giãn dòng và thụt lề đoạn văn
\usepackage{setspace}
\setstretch{1.5}
\usepackage{indentfirst}
\setlength{\parindent}{1.0cm}
\setlength{\parskip}{4pt}

% Các gói đồ họa, bảng biểu và viền trang bìa
\usepackage{graphicx}
\usepackage{xcolor}
\usepackage{tikz}
\usetikzlibrary{calc,shapes.geometric,arrows}
\usepackage{booktabs}
\usepackage{array}
\usepackage{tabularx}
\usepackage{url}
\usepackage[hidelinks]{hyperref}

% Cấu hình Header, Footer và Đánh số trang góc dưới bên phải
\usepackage{fancyhdr}
\pagestyle{fancy}
\fancyhf{}
\fancyfoot[R]{\thepage}
\renewcommand{\headrulewidth}{0pt}
\renewcommand{\footrulewidth}{0pt}

% Cấu hình định dạng tiêu đề
\usepackage{titlesec}
\titleformat{\section}{\fontsize{14}{17}\selectfont\bfseries}{}{0em}{}
\titleformat{\subsection}{\fontsize{13}{16}\selectfont\bfseries}{}{0em}{}
\titleformat{\subsubsection}{\fontsize{13}{16}\selectfont\bfseries\itshape}{}{0em}{}
\titlespacing*{\section}{0pt}{16pt}{7pt}
\titlespacing*{\subsection}{0pt}{12pt}{5pt}
\titlespacing*{\subsubsection}{0pt}{8pt}{3pt}

% Cấu hình mục lục có dấu chấm nối cho tất cả các cấp
\usepackage{tocloft}
\renewcommand{\cftsecleader}{\cftdotfill{\cftdotsep}}
\setlength{\cftsecindent}{0pt}
\setlength{\cftsubsecindent}{15pt}
\setlength{\cftsubsubsecindent}{30pt}

\begin{document}

% 1. Trang bìa chuẩn của Đại học Bách khoa Hà Nội
\begin{titlepage}
    % Khung viền đơn mảnh bám theo lề giấy chuẩn của Đại học Bách khoa Hà Nội
    \begin{tikzpicture}[overlay,remember picture]
        \draw [line width=1.5pt] ($ (current page.north west) + (3cm,-2cm) $) rectangle ($ (current page.south east) + (-2cm,1.5cm) $);
    \end{tikzpicture}
    
    \centering
    
    % Tên Trường & Khoa chủ quản
    {\fontsize{14}{16}\selectfont\textbf{ĐẠI HỌC BÁCH KHOA HÀ NỘI}\\} 
    {\fontsize{12}{14}\selectfont\textbf{KHOA KHOA HỌC VÀ CÔNG NGHỆ GIÁO DỤC}\\} 
    
    \vspace{1.5cm}
    
    % Khối định vị LOGO Trường (tương thích khi có hoặc chưa có tệp ảnh)
    \IfFileExists{hust_logo.png}{%
        \includegraphics[height=4.5cm]{hust_logo.png}%
    }{%
        \fbox{\begin{minipage}[c][4.5cm][c]{4.5cm}\centering\textbf{\color{red!80!black}HUST LOGO}\\\vspace{4pt}\fontsize{9}{11}\selectfont(Đặt file hust\_logo.png tại đây)\end{minipage}}%
    }
    
    \vspace{1.5cm}
    
    % Loại hình tài liệu (Màu đỏ, in hoa đậm)
    {\fontsize{16}{18}\selectfont\textbf{\color{red}TIỂU LUẬN GIỮA KỲ}\\} 
    
    \vspace{0.8cm}
    
    % Tên đề tài nghiên cứu
    {\fontsize{14}{18}\selectfont\textbf{TÌM HIỂU VỀ TẦNG LỚP NHO SĨ\\ Ở VIỆT NAM THẾ KỶ XVIII}\\} 
    
    \vspace{0.4cm}
    \noindent\rule{0.65\textwidth}{0.6pt} % Đường kẻ ngang phân cách tiêu đề
    \vspace{0.6cm}
    
    % Tên Học phần tương ứng
    {\fontsize{12}{14}\selectfont\textit{\textbf{Học phần:} Triết học}\\} 
    
    \vspace{1.8cm}
    
    % Khối thông tin Giảng viên và Học viên thực hiện
    \begin{flushleft}
        \hspace{2.0cm} 
        \begin{tabular}{ll}
            \fontsize{12}{14}\selectfont\textbf{Giảng viên giảng dạy:} & \fontsize{12}{14}\selectfont PGS. TS. Hoàng Thu Hương \\
            \fontsize{12}{14}\selectfont\textbf{Giảng viên nguồn (GVN):} & \fontsize{12}{14}\selectfont Nguyễn Trần Khánh Phương \\
            \fontsize{12}{14}\selectfont\textbf{Học viên thực hiện:} & \fontsize{12}{14}\selectfont Lê Hạnh Nguyên -- 20252233M
        \end{tabular}
    \end{flushleft}
    
    \vfill
    
    % Địa điểm và thời gian xuất bản ở cuối trang
    {\fontsize{12}{14}\selectfont\textbf{Hà Nội -- 2026}\\}
    \vspace*{0.4cm}
    
\end{titlepage}



\newpage
% 2. Trang mục lục (không đánh số trang)
\pagenumbering{gobble}
\renewcommand{\contentsname}{\centering\fontsize{14}{17}\selectfont\textbf{MỤC LỤC}}
{\setstretch{1.2}
\tableofcontents
}
\newpage

% Bắt đầu đánh số trang Ả Rập từ Phần Mở đầu
\pagenumbering{arabic}
\setcounter{page}{1}

% 3. Phần Mở đầu
\section*{PHẦN MỞ ĐẦU}
\addcontentsline{toc}{section}{PHẦN MỞ ĐẦU}

\subsection*{1. Lý do chọn đề tài}
\addcontentsline{toc}{subsection}{1. Lý do chọn đề tài}

Trong hệ thống lý luận của chủ nghĩa duy vật lịch sử, mối quan hệ biện chứng giữa tồn tại xã hội và ý thức xã hội, giữa cơ sở hạ tầng và kiến trúc thượng tầng là nguyên lý nền tảng để nhận thức sự vận động của lịch sử. Triết học Mác -- Lênin chỉ rõ phương thức sản xuất vật chất quyết định đời sống tinh thần của xã hội; ngược lại, ý thức xã hội có tính độc lập tương đối và tác động trở lại thực tiễn thông qua hoạt động của con người. Trong xã hội trung đại Việt Nam, nho sĩ là tầng lớp trực tiếp tiếp nhận, truyền bá và thực thi hệ tư tưởng chính thống, giữ vai trò định hướng chuẩn mực chính trị, văn hóa và đạo đức quốc gia.

Thế kỷ XVIII đánh dấu bước ngoặt lớn của chế độ phong kiến Việt Nam khi thiết chế quân chủ lâm vào khủng hoảng ở cả Đàng Ngoài lẫn Đàng Trong, kéo theo sự suy giảm uy tín của Nho giáo chính thống vốn xác lập vị thế độc tôn từ thời Lê sơ. Cùng thời điểm đó, phong trào nông dân khởi nghĩa bùng nổ rộng khắp và đạt tới đỉnh cao với cuộc khởi nghĩa Tây Sơn do Quang Trung -- Nguyễn Huệ lãnh đạo, xóa bỏ thế phân liệt đất nước và đánh tan các thế lực ngoại xâm Xiêm, Thanh. Trước sự biến đổi căn bản của tồn tại xã hội, tầng lớp nho sĩ không còn giữ được sự thuần nhất về tư tưởng mà phân hóa mạnh mẽ về thái độ chính trị, đồng thời chuyển dịch từ lối học kinh viện Tống Nho sang khuynh hướng thực học, khảo chứng bách khoa và tinh thần nhân đạo.

Từ thực tế lịch sử đó, việc nghiên cứu đề tài ``Tìm hiểu về tầng lớp nho sĩ ở Việt Nam thế kỷ XVIII'' đặt ra yêu cầu thiết thực trên cả hai bình diện lý luận và thực tiễn. Về mặt lý luận, đề tài làm sáng tỏ cơ chế tác động của điều kiện kinh tế -- xã hội đối với sự biến đổi hệ tư tưởng, qua đó nhận diện vị trí trung gian và quy luật phân hóa của trí thức phong kiến trong các giai đoạn chuyển giao lịch sử. Về mặt thực tiễn, những kinh nghiệm lịch sử của kẻ sĩ thế kỷ XVIII gợi mở nhiều giá trị tham chiếu cho công tác xây dựng đội ngũ trí thức Việt Nam hiện nay: gắn kết lý luận với thực tiễn sản xuất, khắc phục bệnh chuộng hư danh và bằng cấp hình thức, bồi dưỡng bản lĩnh chính trị của người trí thức và hoàn thiện chính sách trọng dụng hiền tài phục vụ phát triển đất nước.

\subsection*{2. Tổng quan tình hình nghiên cứu}
\addcontentsline{toc}{subsection}{2. Tổng quan tình hình nghiên cứu}

Diện mạo lịch sử và tư tưởng của tầng lớp nho sĩ Việt Nam thế kỷ XVIII đã được giới nghiên cứu triết học, sử học và văn hóa học tiếp cận dưới ba nhóm công trình chính:

Thứ nhất là nhóm công trình nghiên cứu lịch sử tư tưởng và triết học Việt Nam. Tiêu biểu có bộ sách \textit{Lịch sử tư tưởng Việt Nam} (Tập 1, 1993) do Giáo sư, Tiến sĩ Nguyễn Tài Thư chủ biên đã khảo sát quá trình suy thoái của Nho giáo chính thống từ thế kỷ XVI đến thế kỷ XVIII và sự nảy nở của các yếu tố thực học. Trước đó, công trình \textit{Sự phát triển của tư tưởng ở Việt Nam từ thế kỷ XIX đến Cách mạng Tháng Tám} (Tập 1, 1973) của Giáo sư Trần Văn Giàu cũng phân tích cuộc khủng hoảng ý thức hệ phong kiến cuối thế kỷ XVIII như một tiền đề tư tưởng dẫn tới những biến chuyển ở giai đoạn cận đại.

Thứ hai là nhóm công trình lịch sử chế độ phong kiến và văn hóa giáo dục trung đại. Sách \textit{Lịch sử chế độ phong kiến Việt Nam} (Tập III, 1965) của các tác giả Phan Huy Lê, Chu Thiên, Vương Hoàng Tuyên, Đinh Xuân Lâm đã dựng lại bối cảnh mâu thuẫn xã hội, sự rệu rã của bộ máy quan trường và phong trào nông dân thế kỷ XVIII. Ở khía cạnh chế độ khoa cử, \textit{Việt Nam văn hóa sử cương} của Đào Duy Anh và \textit{Đại cương lịch sử văn hóa Việt Nam} (Tập 4, 1998) của Nguyễn Khắc Thuần cung cấp nhiều tư liệu xác thực về quy chế trường thi cùng những tiêu cực thời Lê mạt.

Thứ ba là nhóm chuyên khảo và bài báo khoa học về từng danh nhân, trào lưu học thuật cụ thể. Bài viết của Nguyễn Thị Hồng Phương trên \textit{Tạp chí Nghiên cứu Triết học} (2026) đã phân tích nhân sinh quan nhập thế gắn với quan niệm Lý -- Khí của Lê Quý Đôn; đồng thời bản dịch và khảo cứu tác phẩm \textit{Trúc Lâm tông chỉ nguyên thanh} (2001) làm rõ xu hướng dung hợp Tam giáo của Ngô Thì Nhậm. Cùng với các khảo cứu văn bản học của Nguyễn Kim Sơn (2004) về khuynh hướng khảo chứng học, các nghiên cứu của Đặng Việt Bích (2015) về thái độ của kẻ sĩ Bắc Hà đối với vua Quang Trung, của Quách Hiền (2018) về mạng lưới tư học làng xã, của Cao Văn Thức (2016) về học vị Sinh đồ và chuyên đề tư liệu về danh sư Võ Trường Toản (2020) đã bổ sung nguồn sử liệu phong phú trên cả hai miền Nam -- Bắc.

Kế thừa kết quả của các công trình đi trước, tiểu luận hệ thống hóa bức tranh toàn cảnh về tầng lớp nho sĩ thế kỷ XVIII dưới phương pháp luận duy vật biện chứng và duy vật lịch sử, gắn kết sự phân hóa chính trị với bước chuyển học thuật để rút ra bài học đối với đội ngũ trí thức đương đại.

\subsection*{3. Mục tiêu và nhiệm vụ nghiên cứu}
\addcontentsline{toc}{subsection}{3. Mục tiêu và nhiệm vụ nghiên cứu}

Mục tiêu tổng quát của đề tài là vận dụng lý luận triết học Mác -- Lênin để làm rõ nguồn gốc xã hội, thực trạng phân hóa chính trị và sự vận động tư tưởng -- học thuật của tầng lớp nho sĩ Việt Nam thế kỷ XVIII; trên cơ sở đó đánh giá khách quan những đóng góp, hạn chế lịch sử và rút ra bài học đối với việc phát huy vai trò đội ngũ trí thức Việt Nam hiện nay.

Để thực hiện mục tiêu trên, đề tài giải quyết ba nhiệm vụ cụ thể tương ứng với ba chương:

\textit{Nhiệm vụ thứ nhất,} xác lập cơ sở lý luận duy vật lịch sử về ý thức xã hội và đặc trưng của tầng lớp trí thức; phân tích vị thế nho sĩ trước thế kỷ XVIII cùng những biến động kinh tế, chính trị, khoa cử thế kỷ XVIII tác động đến tầng lớp này.

\textit{Nhiệm vụ thứ hai,} phân tích các khuynh hướng ứng xử chính trị của nho sĩ trước biến cố thời cuộc và làm sáng tỏ những chuyển biến về thế giới quan triết học, phương pháp khảo chứng, tinh thần dung hợp tôn giáo và giá trị nhân văn trong sáng tác của họ.

\textit{Nhiệm vụ thứ ba,} đánh giá giá trị tiến bộ và giới hạn lịch sử của nho sĩ thế kỷ XVIII theo nguyên tắc lịch sử -- cụ thể, từ đó liên hệ rút ra các bài học thực tiễn trong công tác xây dựng đội ngũ trí thức của Đảng và Nhà nước hiện nay.

\subsection*{4. Đối tượng và phạm vi nghiên cứu}
\addcontentsline{toc}{subsection}{4. Đối tượng và phạm vi nghiên cứu}

Đối tượng nghiên cứu của đề tài là tầng lớp nho sĩ Việt Nam thế kỷ XVIII, tập trung vào sự phân hóa lập trường chính trị -- xã hội và sự chuyển biến trong thế giới quan triết học, nhân sinh quan, khuynh hướng học thuật dưới sự quy định của hoàn cảnh lịch sử đương thời.

Về phạm vi không gian, đề tài khảo sát hoạt động và tư tưởng của nho sĩ trên cả hai vùng lãnh thổ Đàng Ngoài và Đàng Trong của nước Đại Việt. Về phạm vi thời gian, nội dung nghiên cứu giới hạn trọng tâm trong thế kỷ XVIII (từ thời điểm khủng hoảng của chính quyền Lê -- Trịnh, chúa Nguyễn đến giai đoạn triều đại Tây Sơn), có đối chiếu với tiền đề Nho giáo các thế kỷ trước và liên hệ với thực tiễn Việt Nam hiện nay.

\subsection*{5. Phương pháp nghiên cứu}
\addcontentsline{toc}{subsection}{5. Phương pháp nghiên cứu}

Đề tài này được thực hiện trên cơ sở phương pháp luận của chủ nghĩa duy vật biện chứng và chủ nghĩa duy vật lịch sử. Ngoài ra, đề tài còn sử dụng các phương pháp khác như:

Phương pháp thống nhất lịch sử và logic giúp tái hiện diễn biến lịch sử thế kỷ XVIII theo trình tự thời gian, qua đó nắm bắt quy luật vận động nội tại của tư tưởng nho sĩ mà không sa vào lối liệt kê sự kiện thuần túy.

Phương pháp phân tích và tổng hợp được áp dụng để bóc tách từng thái độ xuất xử, từng luận điểm triết học của các học giả tiêu biểu, từ đó khái quát thành diện mạo chung của trí thức thời kỳ này.

Phương pháp so sánh, đối chiếu giúp nhận diện sự khác biệt giữa bộ phận nho sĩ bảo thủ và bộ phận nho sĩ cấp tiến, giữa đặc trưng kẻ sĩ Bắc Hà với sĩ phu Gia Định, cũng như giữa Lý học Tống Nho kinh viện với trào lưu thực học.

Phương pháp tiếp cận hệ thống đặt tầng lớp nho sĩ trong mối liên hệ chỉnh thể với phương thức sản xuất phong kiến, thiết chế khoa cử và phong trào nông dân khởi nghĩa thế kỷ XVIII.

\subsection*{6. Kết cấu của đề tài}
\addcontentsline{toc}{subsection}{6. Kết cấu của đề tài}

Ngoài phần mở đầu, kết luận, tài liệu tham khảo, nội dung của đề tài này bao gồm 03 Chương, 06 Tiết. Mạch nghiên cứu đi từ nền tảng lý luận và điều kiện lịch sử khách quan ở Chương 1, sang phân tích thực trạng phân hóa chính trị cùng bước chuyển tư tưởng ở Chương 2, trước khi tổng kết giá trị, hạn chế và đúc rút bài học thực tiễn tại Chương 3.

\newpage

% 4. Chương 1: Cơ sở lý luận và bối cảnh lịch sử - xã hội
\section*{CHƯƠNG 1: CƠ SỞ LÝ LUẬN VÀ BỐI CẢNH LỊCH SỬ -- XÃ HỘI TÁC ĐỘNG ĐẾN TẦNG LỚP NHO SĨ VIỆT NAM THẾ KỶ XVIII}
\addcontentsline{toc}{section}{CHƯƠNG 1: CƠ SỞ LÝ LUẬN VÀ BỐI CẢNH LỊCH SỬ -- XÃ HỘI TÁC ĐỘNG ĐẾN TẦNG LỚP NHO SĨ VIỆT NAM THẾ KỶ XVIII}

\subsection*{1.1. Quan điểm triết học Mác -- Lênin về ý thức xã hội và khái lược diện mạo nho sĩ trước thế kỷ XVIII}
\addcontentsline{toc}{subsection}{1.1. Quan điểm triết học Mác -- Lênin về ý thức xã hội và khái lược diện mạo nho sĩ trước thế kỷ XVIII}

\subsubsection*{1.1.1. Nguyên lý duy vật lịch sử về tồn tại xã hội, ý thức xã hội và đặc điểm tầng lớp trí thức}
\addcontentsline{toc}{subsubsection}{1.1.1. Nguyên lý duy vật lịch sử về tồn tại xã hội, ý thức xã hội và đặc điểm tầng lớp trí thức}

Chủ nghĩa duy vật lịch sử giải thích sự vận động của xã hội loài người xuất phát từ mối quan hệ biện chứng giữa tồn tại xã hội và ý thức xã hội. Tồn tại xã hội là toàn bộ phương diện sinh hoạt vật chất và các điều kiện sinh hoạt vật chất của xã hội, bao gồm phương thức sản xuất vật chất, điều kiện tự nhiên -- địa lý và dân cư, trong đó phương thức sản xuất giữ vai trò quyết định. Ý thức xã hội là phương diện sinh hoạt tinh thần của xã hội, nảy sinh từ tồn tại xã hội và phản ánh tồn tại xã hội trong những giai đoạn lịch sử nhất định thông qua các hình thái như ý thức chính trị, pháp quyền, đạo đức, tôn giáo, nghệ thuật, khoa học và triết học.

Các Mác và Phriđrích Ăngghen khẳng định tồn tại xã hội quyết định nội dung, tính chất và xu hướng biến đổi của ý thức xã hội; khi phương thức sản xuất và các quan hệ kinh tế thay đổi thì những tư tưởng, quan điểm chính trị, đạo đức sớm muộn cũng biến đổi theo. Song học thuyết mácxít đồng thời chỉ rõ ý thức xã hội có tính độc lập tương đối, thể hiện trên bốn khía cạnh: ý thức xã hội thường lạc hậu hơn tồn tại xã hội do sức ỳ của tập quán và sự níu kéo từ lợi ích của các giai cấp thống trị cũ; một số tư tưởng khoa học, tiến bộ có thể vượt trước, dự báo xu hướng vận động của hiện thực; sự phát triển của ý thức xã hội luôn mang tính kế thừa lịch sử; và đặc biệt, các hình thái ý thức xã hội có sự tương tác lẫn nhau cũng như tác động trở lại tồn tại xã hội theo hai chiều hướng thúc đẩy hoặc kìm hãm, tùy thuộc vào mức độ phù hợp của hệ tư tưởng đó với quy luật khách quan.

Trong kết cấu xã hội -- giai cấp, triết học Mác -- Lênin xác định trí thức là một tầng lớp xã hội đặc thù chuyên lao động trí óc phức tạp, trực tiếp sáng tạo, bảo lưu và truyền bá tri thức văn hóa, tư tưởng. Trí thức không cấu thành một giai cấp kinh tế độc lập vì họ không nắm giữ tư liệu sản xuất riêng biệt và không đại diện cho một phương thức sản xuất riêng. Chính vị trí trung gian về mặt kinh tế khiến lập trường chính trị và thế giới quan của trí thức luôn chịu sự quy định bởi các giai cấp cơ bản trong xã hội. Khi chế độ xã hội ổn định và giai cấp cầm quyền còn đóng vai trò tiến bộ, trí thức là lực lượng nòng cốt củng cố thiết chế và hệ tư tưởng chính thống. Ngược lại, khi phương thức sản xuất cũ lâm vào bế tắc và mâu thuẫn giai cấp gay gắt, tầng lớp trí thức tất yếu mất đi tính thống nhất nội tại, buộc phải lựa chọn lại chỗ đứng chính trị và định hướng học thuật trước những đòi hỏi mới của lịch sử.

\subsubsection*{1.1.2. Khái niệm nho sĩ và vị thế của tầng lớp nho sĩ Việt Nam trong lịch sử trung đại trước thế kỷ XVIII}
\addcontentsline{toc}{subsubsection}{1.1.2. Khái niệm nho sĩ và vị thế của tầng lớp nho sĩ Việt Nam trong lịch sử trung đại trước thế kỷ XVIII}

Trong xã hội phong kiến Việt Nam, nho sĩ (hay kẻ sĩ, sĩ phu Nho học) là tầng lớp trí thức được đào tạo bằng hệ thống kinh điển Nho gia (\textit{Tứ thư}, \textit{Ngũ kinh}), tham gia khoa cử để bước vào bộ máy quan liêu hoặc lui về làng xã dạy học, bốc thuốc, giữ gìn hương ước. Trong trật tự phân tầng ``Sĩ, Nông, Công, Thương'', kẻ sĩ đứng ở vị trí đầu bảng nhờ uy tín học vấn và vai trò biểu trưng đạo lý. Nhân sinh quan chính trị của họ xoay quanh lộ trình ``tu thân, tề gia, trị quốc, bình thiên hạ'' và nguyên tắc ứng xử ``xuất -- xử'': gặp thời thịnh trị, vua sáng tôi hiền thì xuất thế hành đạo giúp nước; gặp buổi hôn quân bạo chúa, kỷ cương đảo lộn thì ẩn cư giữ trọn khí tiết người quân tử.

Nhìn lại chặng đường trước thế kỷ XVIII, vị thế của tầng lớp nho sĩ gắn liền với quá trình xác lập và biến đổi của Nho giáo tại Đại Việt. Dưới triều Lý -- Trần (thế kỷ XI -- XIV), Nho giáo từng bước đi vào thiết chế giáo dục và hành chính qua việc xây dựng Văn Miếu (1070), mở khoa thi Minh kinh bác học đầu tiên (1075) và lập Quốc Tử Giám (1076). Tuy vậy, trong suốt thời kỳ này, Nho giáo chưa giữ địa vị độc tôn mà dung hợp hài hòa với Phật giáo và Đạo giáo trên nền tảng ``Tam giáo đồng nguyên'', tạo nên lớp nho sĩ vừa mẫn cán việc triều chính vừa cởi mở, phóng khoáng về tư duy.

Sang thời Lê sơ thế kỷ XV, đỉnh cao là triều vua Lê Thánh Tông, nhà nước quân chủ trung ương tập quyền đưa Nho giáo (theo hệ phái Tống Nho của Trình -- Chu) lên địa vị độc tôn làm bệ đỡ tư tưởng chính thống. Bài ký trên bia Tiến sĩ khoa Nhâm Tuất dựng tại Văn Miếu năm 1484 do Thân Nhân Trung soạn đã đúc kết quan điểm cốt lõi của thời đại: ``Hiền tài là nguyên khí của quốc gia, nguyên khí thịnh thì thế nước mạnh rồi lên cao, nguyên khí suy thì thế nước yếu rồi xuống thấp''. Nho sĩ khoa bảng trở thành rường cột trực tiếp vận hành nhà nước pháp chế Hồng Đức, đại diện cho kỷ cương phong kiến đang ở giai đoạn hưng thịnh.

Đến thế kỷ XVI -- XVII, sự suy yếu của nhà Lê sơ đẩy đất nước vào cục diện phân tranh kéo dài giữa Nam -- Bắc triều (Lê -- Mạc) và tiếp đó là Trịnh -- Nguyễn phân tranh chia cắt hai miền Đàng Ngoài, Đàng Trong. Thực tế vua Lê chỉ còn hư vị, quyền bính thuộc về các chúa đã làm rạn nứt quan niệm ``chính danh'' và thuyết ``thiên mệnh'' Nho gia. Một số bậc đại nho tiêu biểu như Trạng Trình Nguyễn Bỉnh Khiêm đã chủ động cáo quan về dựng am Bạch Vân dạy học, bày tỏ thái độ phản tỉnh trước các cuộc nội chiến huynh đệ tương tàn. Dù vậy, trong hai thế kỷ này, hệ chuẩn luân lý tam cương ngũ thường về cơ bản vẫn kiểm soát được trật tự xã hội truyền thống.

\subsection*{1.2. Điều kiện kinh tế, chính trị, văn hóa thế kỷ XVIII và những biến động tác động đến tầng lớp nho sĩ}
\addcontentsline{toc}{subsection}{1.2. Điều kiện kinh tế, chính trị, văn hóa thế kỷ XVIII và những biến động tác động đến tầng lớp nho sĩ}

\subsubsection*{1.2.1. Sự khủng hoảng toàn diện của thể chế phong kiến Đàng Ngoài -- Đàng Trong và sự suy đồi của khoa cử}
\addcontentsline{toc}{subsubsection}{1.2.1. Sự khủng hoảng toàn diện của thể chế phong kiến Đàng Ngoài -- Đàng Trong và sự suy đồi của khoa cử}

Bước sang thế kỷ XVIII, chế độ phong kiến ở cả hai miền Nam -- Bắc đồng loạt lâm vào khủng hoảng cấu trúc. Tại Đàng Ngoài, thể chế lưỡng đầu Vua Lê -- Chúa Trịnh bộc lộ sự mục ruỗng từ thượng tầng: danh nghĩa ``phù Lê'' chỉ còn là vỏ bọc cho nạn chuyên quyền, kiêu binh nổi loạn và những cuộc thanh trừng phe phái trong phủ chúa. Ở địa phương, quan lại cấu kết với địa chủ cường hào ráo riết kiêm tính ruộng đất công tư, đẩy nông dân vào cảnh mất tư liệu sản xuất, bần cùng hóa và lưu tán hàng loạt giữa những nạn đói lớn. Tại Đàng Trong, sau giai đoạn đầu khai phá năng động, chính quyền chúa Nguyễn nửa sau thế kỷ XVIII cũng suy đồi nhanh chóng, điển hình dưới thời chúa Nguyễn Phúc Thuần khi quyền thần Trương Phúc Loan thao túng triều chính, mua quan bán tước và đặt ra chế độ thuế khóa nặng nề.

Trên phương diện kinh tế -- xã hội, bất chấp chính sách ``trọng nông ức thương'' của nhà nước phong kiến, kinh tế hàng hóa và luồng giao thương đô thị ở thế kỷ XVIII (tại Kẻ Chợ, Phố Hiến, Hội An, Cù lao Phố, Gia Định) vẫn có bước phát triển khởi sắc. Sự thâm nhập của quan hệ tiền tệ vào đời sống xã hội một mặt kích thích nhu cầu tri thức thực tiễn và ý thức thị dân, mặt khác khi kết hợp với bộ máy quan liêu suy đồi lại làm trầm trọng thêm tệ nạn thương mại hóa quyền lực và học vị.

Biểu hiện trực diện nhất tác động đến tâm thức nho sĩ là sự suy thoái của nền giáo dục và khoa cử quan phương. Quốc Tử Giám cùng hệ thống trường học phủ, huyện sa sút, kỷ cương trường thi buông lỏng với nạn mang sách, mướn người làm bài thay. Đỉnh điểm diễn ra vào năm Canh Ngọ 1750 dưới thời vua Lê Hiển Tông -- chúa Trịnh Doanh: do ngân khố kiệt quệ vì chiến phí đàn áp khởi nghĩa nông dân, triều đình chấp thuận đề xuất của Tham tụng Đỗ Thế Giai, cho phép người dự thi chỉ cần nộp ba quan ``tiền thông kinh'' là được miễn kỳ khảo hạch ở địa phương để vào thẳng trường thi Hương.

Việc xóa bỏ khâu sát hạch năng lực đầu vào khiến trường thi năm đó hỗn loạn chưa từng thấy; người buôn bán, kẻ giàu có không thông chữ nghĩa cũng ùn ùn nộp tiền vào trường chép mướn, dẫn đến cảnh đỗ đạt tràn lan và sinh ra câu châm biếm dân gian ``Sinh đồ ba quan''. Khi con đường khoa bảng -- vốn được kẻ sĩ tôn thờ là thước đo tài đức công minh -- bị đồng tiền thao túng công khai, niềm tin của giới nho sĩ chân chính vào tính chính đáng của triều đình và trật tự lễ giáo chính thống đã đổ vỡ nghiêm trọng.

\subsubsection*{1.2.2. Cao trào khởi nghĩa nông dân với đỉnh cao phong trào Tây Sơn và yêu cầu thời cuộc đối với kẻ sĩ}
\addcontentsline{toc}{subsubsection}{1.2.2. Cao trào khởi nghĩa nông dân với đỉnh cao phong trào Tây Sơn và yêu cầu thời cuộc đối với kẻ sĩ}

Áp bức giai cấp nặng nề đã châm ngòi cho các cuộc khởi nghĩa nông dân nổ ra liên tục suốt thế kỷ XVIII, làm lung lay tận gốc nền thống trị phong kiến. Ở Đàng Ngoài, các cuộc nổi dậy quy mô lớn của Nguyễn Danh Phương, Nguyễn Tuyển, Nguyễn Cừ, Hoàng Công Chất và đặc biệt là cuộc khởi nghĩa kéo dài hơn một thập niên của Nguyễn Hữu Cầu đã uy hiếp trực tiếp kinh thành Thăng Long, buộc giới sĩ phu phải chứng kiến sức mạnh phản kháng của tầng lớp nông dân vốn bị coi thường dưới đáy xã hội.

Đến năm 1771, cuộc khởi nghĩa Tây Sơn do ba anh em Nguyễn Nhạc, Nguyễn Lữ, Nguyễn Huệ phát động từ ấp Tây Sơn (Bình Định) đã chuyển hóa từ phong trào đấu tranh giai cấp ở địa phương thành cao trào thống nhất đất nước và bảo vệ độc lập dân tộc. Nghĩa quân Tây Sơn lần lượt lật đổ chính quyền chúa Nguyễn ở Đàng Trong (1777), đánh tan 5 vạn quân xâm lược Xiêm tại trận Rạch Gầm -- Xoài Mút (1785), rồi tiến ra Bắc Hà tiêu diệt thế lực chúa Trịnh (1786), chấm dứt hơn hai thế kỷ chia cắt sông Gianh. Đỉnh cao là mùa Xuân năm Kỷ Dậu 1789, dưới sự thống lĩnh của Hoàng đế Quang Trung -- Nguyễn Huệ, đại quân Tây Sơn đã hành quân thần tốc quét sạch 29 vạn quân Mãn Thanh do vua Lê Chiêu Thống rước về giày xéo đất nước, giữ vững nền độc lập tự chủ của Đại Việt.

Sự sụp đổ của các tập đoàn phong kiến Lê -- Trịnh, chúa Nguyễn cùng sự xác lập vương triều mới xuất thân từ áo vải nông dân đã tạo nên cú chấn động lớn đối với toàn bộ tầng lớp nho sĩ. Hệ chuẩn luân lý trung quân cứng nhắc của Tống Nho không còn chỗ dựa thực tế khi chính ông vua cuối cùng của triều Lê đã phản bội Tổ quốc, còn sứ mệnh cứu nước lại thuộc về những người anh hùng nông dân. Thực tại khách quan ấy trở thành phép thử lịch sử khắc nghiệt, thúc đẩy quá trình phân hóa chính trị và tái cấu trúc tư duy học thuật của nho sĩ Việt Nam trong suốt thế kỷ XVIII.

\newpage

% 5. Chương 2: Thực trạng phân hóa và chuyển biến tư tưởng
\section*{CHƯƠNG 2: THỰC TRẠNG PHÂN HÓA LẬP TRƯỜNG VÀ NHỮNG CHUYỂN BIẾN TƯ TƯỞNG CỦA TẦNG LỚP NHO SĨ THẾ KỶ XVIII}
\addcontentsline{toc}{section}{CHƯƠNG 2: THỰC TRẠNG PHÂN HÓA LẬP TRƯỜNG VÀ NHỮNG CHUYỂN BIẾN TƯ TƯỞNG CỦA TẦNG LỚP NHO SĨ THẾ KỶ XVIII}

\subsection*{2.1. Sự phân hóa thái độ chính trị -- xã hội của tầng lớp nho sĩ trước các biến cố lịch sử}
\addcontentsline{toc}{subsection}{2.1. Sự phân hóa thái độ chính trị -- xã hội của tầng lớp nho sĩ trước các biến cố lịch sử}

\subsubsection*{2.1.1. Xu hướng bảo thủ giữ chữ trung quân mù quáng và xu hướng bất mãn ẩn dật thúc đẩy phong trào tư học}
\addcontentsline{toc}{subsubsection}{2.1.1. Xu hướng bảo thủ giữ chữ trung quân mù quáng và xu hướng bất mãn ẩn dật thúc đẩy phong trào tư học}

Trước sự cáo chung của chính quyền Lê -- Trịnh và chúa Nguyễn, mâu thuẫn giữa tín điều ``trung quân'' với một dòng họ và nghĩa vụ đối với vận mệnh dân tộc đã chia rẽ tầng lớp nho sĩ thành những thái độ ứng xử trái ngược nhau.

Ở cực bảo thủ nhất là bộ phận nho sĩ chịu sự chi phối nặng nề của giáo điều ``Trung thần bất sự nhị quân'' (Bề tôi trung không thờ hai chúa) và thuyết ``Chính danh'' Tống Nho. Đối với họ, ngôi vị nhà Lê là thiên mệnh bất biến, còn mọi cuộc nổi dậy của nông dân đều bị quy thành hành vi tiếm nghịch. Khi trật tự cũ sụp đổ trước sức tiến công của quân Tây Sơn, nhiều người rơi vào bế tắc tư tưởng và chọn cách giữ tiết cực đoan: người thì tự vẫn tuẫn tiết theo triều cũ như Lý Trần Quán (tự chôn sống năm 1786), Trần Công Xán (chịu chết khi đi sứ Phú Xuân năm 1787); người thì bôn tẩu cần vương bất thành rồi lui về ẩn cư, cự tuyệt hợp tác với triều đại mới đến trọn đời như Trần Danh Án, Lê Duy Đản.

Nghiêm trọng hơn, sự mù quáng về ý thức hệ đã đẩy nhóm quan lại nho sĩ như Lê Quynh bước qua ranh giới phản bội dân tộc. Nhân danh khôi phục cơ nghiệp nhà Lê, nhóm này đã hộ tống vua Lê Chiêu Thống sang cầu viện triều đình Mãn Thanh, trực tiếp dẫn đường cho quân xâm lược của Tôn Sĩ Nghị tràn vào giày xéo đất nước. Dưới góc độ duy vật lịch sử, đây là bi kịch tất yếu khi lòng trung thành giai cấp hẹp hòi bị đặt cao hơn chủ quyền quốc gia và quyền sống của nhân dân.

Khác với xu hướng bảo thủ trên, một bộ phận đông đảo nho sĩ có lương tri lại chọn con đường thoái lui khỏi chốn quan trường. Họ chán ghét sự mục nát, phe phái trong phủ chúa cung vua, song do giới hạn giai cấp nên vẫn dè dặt, hoài nghi trước phong trào nông dân Tây Sơn. Vận dụng triết lý ``Độc thiện kỳ thân'' (giữ mình trong sạch giữa thời loạn) và ``Dĩ thân tàng'' (ẩn mình chờ thời), nhiều danh nho từ quan trở về quê nhà như Bùi Huy Bích, Phạm Quý Thích, hay La Sơn Phu Tử Nguyễn Thiếp trong giai đoạn ẩn cư tại trại Bùi Phong trên dãy núi Thiên Nhẫn.

Sự thoái lui chính trị của các bậc túc nho này lại mang đến một hệ quả văn hóa tích cực là thúc đẩy mạng lưới tư học phát triển mạnh mẽ. Giữa lúc Quốc Tử Giám và trường học công lập suy tàn, các lớp học tư gia của Bùi Huy Bích ở Định Công hay của Nguyễn Thiếp ở xứ Nghệ đã thu hút hàng trăm môn sinh đến thụ giáo. Tại đây, người thầy không chỉ truyền dạy kinh sử mà còn đề cao đạo làm người, giữ gìn sĩ khí trước tệ nạn mua danh bán tước nơi trường thi. Phong trào tư học thế kỷ XVIII nhờ đó đã bảo toàn mạch nguồn văn hiến và lan tỏa tri thức xuống tận cấp cơ sở làng xã.

\subsubsection*{2.1.2. Xu hướng tiến bộ dấn thân cùng phong trào Tây Sơn và nét đặc thù khai phóng của nho sĩ Đàng Trong}
\addcontentsline{toc}{subsubsection}{2.1.2. Xu hướng tiến bộ dấn thân cùng phong trào Tây Sơn và nét đặc thù khai phóng của nho sĩ Đàng Trong}

Vượt lên trên tâm lý hoài cổ hay ẩn dật, bộ phận nho sĩ ưu tú nhất của Bắc Hà đã dứt khoát rũ bỏ định kiến giai cấp để đồng hành cùng phong trào Tây Sơn. Nhận thức của họ bắt nguồn từ sự kế thừa tư tưởng ``Dân vi bản'' (Dân là gốc của nước): triều đại nào bảo vệ được độc lập dân tộc và đem lại cuộc sống yên ổn cho trăm họ thì triều đại đó nắm giữ tính chính thống lịch sử. Đáp lại tấm lòng trân trọng hiền tài của Hoàng đế Quang Trung qua bản \textit{Chiếu cầu hiền} (do Ngô Thì Nhậm chấp bút) và tư tưởng cải cách trong \textit{Chiếu lập học} -- ``Xây dựng đất nước lấy việc dạy học làm đầu, tìm lẽ trị bình lấy việc tuyển nhân tài làm gấp'' -- nhiều sĩ phu tài danh đã đem trí tuệ phụng sự vương triều mới.

Tiêu biểu nhất cho xu hướng nhập thế này là Ngô Thì Nhậm (1746--1803). Mùa đông năm 1788, trước thế giặc Mãn Thanh ồ ạt tiến công, chính ông đã đề xuất kế sách rút quân chiến lược khỏi Thăng Long về lập phòng tuyến liên hoàn Tam Điệp -- Biện Sơn nhằm bảo toàn chủ lực, làm kiêu lòng giặc, tạo bàn đạp cho cuộc hành quân thần tốc đại phá quân Thanh mùa Xuân năm Kỷ Dậu 1789. Sau ngày chiến thắng, được vua Quang Trung giao trọng trách bang giao, Ngô Thì Nhậm đã vận dụng sách lược ngoại giao vừa kiên quyết giữ vững chủ quyền, vừa mềm dẻo khéo léo, buộc triều đình Càn Long phải bãi bỏ ý định phục thù và chính thức công nhận vương triều Tây Sơn.

Sát cánh cùng Ngô Thì Nhậm là những bậc danh nho đức độ và thực tài. La Sơn Phu Tử Nguyễn Thiếp (1723--1804), người từng nhiều lần khước từ lời vời của chúa Trịnh, đã cảm phục tầm vóc anh hùng và sự cầu thị của vua Quang Trung mà nhận chức Viện trưởng Viện Sùng Chính, chủ trì dịch kinh sách chữ Hán sang chữ Nôm để cải cách giáo dục và tham mưu quy hoạch Phượng Hoàng Trung Đô tại Nghệ An. Trong khi đó, Phan Huy Ích (1751--1822) đảm nhận soạn thảo văn thư bang giao, xây dựng điển chế triều nghi và để lại bộ \textit{Dụ Am văn tập} có giá trị sử liệu cao.

Ở phía Nam, tầng lớp nho sĩ xứ Đàng Trong (đặc biệt tại vùng Gia Định -- Nam Bộ) lại định hình phong cách khai phóng rõ nét. Do sinh trưởng trên vùng đất mới khai phá, nơi hội tụ cư dân đa nguồn gốc cùng hoạt động thương cảng sôi động, kẻ sĩ phương Nam ít bị ràng buộc bởi lễ giáo từ chương khắt khe như Bắc Hà, mà giàu tính thực tiễn, trọng nghĩa khí và gần gũi với dân nghèo.

Gương mặt đại diện cho sĩ phu Nam Bộ thế kỷ XVIII là danh sư Võ Trường Toản (mất năm 1792), bậc ``Sùng Đức xử sĩ'' được người đương thời tôn vinh là ``Sư biểu đất Gia Định''. Giữa cuộc giao tranh quyết liệt giữa Nguyễn Ánh và Tây Sơn, ông giữ thái độ độc lập, không ra làm quan mà mở trường dạy học tại làng Hòa Hưng. Võ Trường Toản đề xướng phương châm giáo dục ``Tri ngôn dưỡng khí'' (hiểu thấu đạo lý để nuôi dưỡng chí khí khẳng khái), chú trọng rèn luyện thực đức hơn lối văn chương cử nghiệp. Trường học của ông đã đào tạo nên lớp trí thức trụ cột của miền Nam, nổi bật là nhóm ``Gia Định tam gia'' gồm Trịnh Hoài Đức, Lê Quang Định và Ngô Nhân Tịnh -- những tác giả đã biên soạn bộ địa chí \textit{Gia Định thành thông chí} và góp phần hun đúc nên ``Hào khí Đồng Nai''.

\subsection*{2.2. Những bước chuyển biến trong thế giới quan triết học, khuynh hướng học thuật và ý thức nhân văn}
\addcontentsline{toc}{subsection}{2.2. Những bước chuyển biến trong thế giới quan triết học, khuynh hướng học thuật và ý thức nhân văn}

\subsubsection*{2.2.1. Bước chuyển từ Lý học Tống Nho kinh viện sang trào lưu thực học và tư duy bách khoa khảo chứng}
\addcontentsline{toc}{subsubsection}{2.2.1. Bước chuyển từ Lý học Tống Nho kinh viện sang trào lưu thực học và tư duy bách khoa khảo chứng}

Suốt từ thế kỷ XV, dòng Tống Nho (phái Trình -- Chu) chi phối học thuật Đại Việt với xu hướng thiên về bình giải kinh viện các phạm trù trừu tượng như Tâm, Tính, Mệnh và quy phạm hóa đạo đức tam cương ngũ thường. Bước sang thế kỷ XVIII, lối học tư biện sáo rỗng ấy tỏ ra bất lực trước hàng loạt vấn đề cấp bách của đất nước như trị thủy, cứu đói, quản lý điền chế và phát triển giao thương. Từ yêu cầu thực tiễn đó, giới nho sĩ tiến bộ đã chuyển hướng mạnh mẽ sang trào lưu thực học và phương pháp khảo chứng bách khoa, đề cao tri thức gắn với ``quốc kế dân sinh''.

Đỉnh cao của khuynh hướng thực học và tư duy triết học thế kỷ XVIII kết tinh ở nhà bác học Lê Quý Đôn (1726--1784) trên bốn phương diện chủ yếu:

Về bản thể luận, trong mục \textit{Lý khí} của bộ sách \textit{Vân đài loại ngữ}, Lê Quý Đôn dùng phạm trù ``Khí'' để cắt nghĩa cội nguồn vật chất của vũ trụ. Ông quan niệm Khí là thực thể nguyên sơ lấp đầy không gian, sự tụ tán và vận động tự thân của âm dương nhị khí tạo thành muôn vật trong trời đất. Luận điểm này mang khuynh hướng duy vật tự phát rõ rệt, đối lập với cách giải thích thế giới bằng ý chí thần linh siêu nhiên.

Về quan hệ giữa Lý và Khí, Lê Quý Đôn trực tiếp phản bác mệnh đề duy tâm khách quan của Chu Hy (vốn đặt Lý có trước và sinh ra Khí). Ông khẳng định ``Lý ở trong Khí'', coi Lý không phải là thực thể tinh thần tồn tại biệt lập bên ngoài vật chất, mà chỉ là quy luật xác định vốn có của từng sự vật hữu hình. Do quy luật nằm ngay trong sự vật, muốn nhận thức đúng đắn thì con người phải khảo sát trực tiếp bản thân thế giới hiện thực.

Về phương pháp nhận thức, Lê Quý Đôn giải thích lại mệnh đề ``cách vật trí tri'' theo hướng thực chứng: nghiên cứu phải dựa trên thu thập chứng cứ, đối chiếu thư tịch cổ kim và khảo sát thực địa. Ông để lại hệ thống công trình bách khoa đồ sộ gồm \textit{Vân đài loại ngữ}, \textit{Kiến văn tiểu lục}, \textit{Đại Việt thông sử} và \textit{Phủ biên tạp lục}. Riêng trong \textit{Phủ biên tạp lục}, những trang ghi chép tỉ mỉ của ông về vị trí địa lý, sản vật cùng hoạt động thực thi chủ quyền của đội Hoàng Sa và đội Bắc Hải đã trở thành nguồn sử liệu gốc xác thực về chủ quyền của Việt Nam đối với hai quần đảo Hoàng Sa, Trường Sa.

Về tư duy kinh tế, Lê Quý Đôn vượt qua định kiến ``trọng nông ức thương'' truyền thống. Qua các khảo cứu trong mục \textit{Điển vựng}, \textit{Phẩm vật} sách \textit{Vân đài loại ngữ} và \textit{Phủ biên tạp lục}, ông chủ trương kết hợp hài hòa giữa phát triển nông nghiệp với mở mang thủ công nghiệp, khai mỏ, lưu thông tiền tệ -- thương mại và bồi dưỡng tầng lớp hiền tài. Tinh thần kinh tế toàn diện đó thường được hậu thế đúc kết qua câu phương ngôn ``Phi nông bất ổn, phi công bất phú, phi thương bất hoạt, phi trí bất hưng'', phản ánh nhận thức biện chứng về bốn trụ cột giữ cho quốc gia ổn định, giàu mạnh và hưng thịnh.

\subsubsection*{2.2.2. Khuynh hướng hòa đồng Tam giáo và sự trỗi dậy của chủ nghĩa nhân đạo trong sáng tác của nho sĩ}
\addcontentsline{toc}{subsubsection}{2.2.2. Khuynh hướng hòa đồng Tam giáo và sự trỗi dậy của chủ nghĩa nhân đạo trong sáng tác của nho sĩ}

Song hành cùng trào lưu thực học là sự phục hưng tinh thần dung hợp Tam giáo (Nho -- Phật -- Đạo). Khi khuôn mẫu giáo điều Nho gia không còn giải đáp thỏa đáng những bi kịch nhân sinh thời loạn, nhiều nho sĩ đã tìm đến triết lý từ bi, vô ngã của Phật giáo và quan niệm thuận theo tự nhiên của Lão -- Trang để cân bằng đời sống tinh thần.

Công trình tiêu biểu nhất của xu hướng này là chuyên luận triết học \textit{Trúc Lâm tông chỉ nguyên thanh} do Ngô Thì Nhậm biên soạn dưới pháp danh Hải Lượng đại thiền sư. Kế thừa tư tưởng nhập thế ``Cư trần lạc đạo'' của Thiền phái Trúc Lâm thời Trần, Ngô Thì Nhậm chủ trương thống nhất Nho với Phật trên cùng một bản thể tâm tính: người trí thức khi xuất thế gánh vác việc quân quốc thì hành xử theo đạo Nho, khi quay về tu dưỡng nội tâm để giữ sự thanh tịnh, sáng suốt thì thấm nhuần đạo Thiền. Cách tiếp cận đó đã phá vỡ ranh giới bài xích tôn giáo hẹp hòi của Tống Nho từng tồn tại suốt nhiều thế kỷ.

Sự cởi mở về thế giới quan đồng thời khơi nguồn cho trào lưu chủ nghĩa nhân đạo phát triển rực rỡ trong văn học chữ Nôm của giới nho sĩ. Văn chương thế kỷ XVIII thoát dần khỏi công thức giáo huấn ``văn dĩ tải đạo'' để hướng thẳng vào số phận con người trần thế, lên án chiến tranh phong kiến và đòi quyền sống, quyền hạnh phúc cá nhân.

Mỗi tác phẩm lớn thời kỳ này đều mang giá trị phê phán hiện thực và cảm quan nhân văn sắc nét. Khúc ngâm \textit{Chinh phụ ngâm khúc} (nguyên tác chữ Hán của Đặng Trần Côn, bản diễn Nôm truyền tụng gắn với Đoàn Thị Điểm và Phan Huy Ích) là tiếng nói xót thương tuổi xuân người phụ nữ bị chia lìa bởi chiến tranh phi nghĩa. Trong khi đó, \textit{Cung oán ngâm khúc} của Nguyễn Gia Thiều phơi bày bi kịch tù túng nơi cung cấm để phê phán lối sống hưởng lạc của bậc đế vương. Ở mạch thơ Nôm truyền tụng, Hồ Xuân Hương dùng tiếng cười trào phúng đả kích trực diện thói đạo đức giả của giới nho sĩ hủ lậu và khẳng định khát vọng bình đẳng của người phụ nữ; còn truyện thơ \textit{Sơ kính tân trang} của Phạm Thái lại cất lên tiếng nói bảo vệ tình yêu tự do trước lễ giáo môn đăng hộ đối khắc nghiệt. Những thành tựu ấy đã tạo tiền đề nghệ thuật và tư tưởng trực tiếp cho sự ra đời kiệt tác \textit{Đoạn trường tân thanh} (Truyện Kiều) của Nguyễn Du vào đầu thế kỷ XIX.

Như vậy, biến động xã hội thế kỷ XVIII đã thúc đẩy tầng lớp nho sĩ Việt Nam thực hiện một cuộc tự vượt bỏ trên cả hai bình diện chính trị và học thuật. Từ sự phân hóa thái độ xuất xử đến sự hình thành trào lưu thực học, dung hợp Tam giáo và chủ nghĩa nhân đạo, diện mạo tinh thần của kẻ sĩ Đại Việt đã đạt tới độ chín muồi mới trong lịch sử tư tưởng dân tộc.

\newpage

% 6. Chương 3: Giá trị, hạn chế lịch sử và bài học đối với đội ngũ trí thức
\section*{CHƯƠNG 3: GIÁ TRỊ, HẠN CHẾ LỊCH SỬ CỦA TẦNG LỚP NHO SĨ THẾ KỶ XVIII VÀ BÀI HỌC ĐỐI VỚI ĐỘI NGŨ TRÍ THỨC VIỆT NAM HIỆN NAY}
\addcontentsline{toc}{section}{CHƯƠNG 3: GIÁ TRỊ, HẠN CHẾ LỊCH SỬ CỦA TẦNG LỚP NHO SĨ THẾ KỶ XVIII VÀ BÀI HỌC ĐỐI VỚI ĐỘI NGŨ TRÍ THỨC VIỆT NAM HIỆN NAY}

\subsection*{3.1. Những đóng góp tiến bộ và hạn chế lịch sử của tầng lớp nho sĩ thế kỷ XVIII dưới góc độ duy vật lịch sử}
\addcontentsline{toc}{subsection}{3.1. Những đóng góp tiến bộ và hạn chế lịch sử của tầng lớp nho sĩ thế kỷ XVIII dưới góc độ duy vật lịch sử}

\subsubsection*{3.1.1. Những giá trị đóng góp về mặt tư tưởng, văn hóa giáo dục và tinh thần độc lập dân tộc}
\addcontentsline{toc}{subsubsection}{3.1.1. Những giá trị đóng góp về mặt tư tưởng, văn hóa giáo dục và tinh thần độc lập dân tộc}

Vận dụng nguyên tắc toàn diện và lịch sử -- cụ thể của chủ nghĩa duy vật lịch sử, có thể khẳng định sự vận động của tầng lớp nho sĩ thế kỷ XVIII đã kết tinh thành bốn giá trị đóng góp nổi bật đối với tiến trình phát triển của dân tộc:

Thứ nhất, trên bình diện tư tưởng triết học và phương pháp luận khoa học, các học giả tiêu biểu thế kỷ XVIII đã bước đầu giải phóng tư duy lý luận Việt Nam khỏi sự ràng buộc của chủ nghĩa duy tâm khách quan và lối học tầm chương trích cú Tống Nho. Việc đề cao bản thể luận vật chất tự phát về ``Khí'', xác lập mệnh đề ``Lý ở trong Khí'' cùng phương pháp khảo chứng thực địa đã đặt nền móng cho lối tư duy trọng thực chứng, hướng tri thức vào giải quyết các vấn đề kinh tế, điển chế, địa lý và xác lập chứng cứ chủ quyền lãnh thổ, biển đảo quốc gia.

Thứ hai, trên bình diện văn hóa -- giáo dục, tầng lớp nho sĩ đã thực hiện thành công quá trình xã hội hóa và bản địa hóa tri thức trong bối cảnh thiết chế giáo dục quan phương suy sụp. Mạng lưới trường tư thục rộng khắp ở nông thôn Bắc Hà và vùng đất mới Gia Định, kết hợp với nỗ lực đưa chữ Nôm vào dịch thuật kinh điển và văn thư hành chính dưới triều Tây Sơn, đã phá vỡ thế độc quyền chữ nghĩa của cung đình, giữ vững truyền thống hiếu học và đạo lý dân tộc ngay từ nền tảng làng xã.

Thứ ba, trên bình diện chính trị và bảo vệ Tổ quốc, bộ phận nho sĩ cấp tiến đã chứng minh vai trò tích cực của ý thức xã hội tiên tiến khi tác động trở lại thực tiễn lịch sử. Bằng việc vượt qua quan niệm trung quân đóng khung trong một dòng họ phong kiến để lấy lợi ích dân tộc và tư tưởng ``dân vi bản'' làm thước đo chính nghĩa, trí tuệ quân sự -- ngoại giao của kẻ sĩ đã kết hợp nhuần nhuyễn với sức mạnh quật khởi của phong trào nông dân, góp phần quyết định vào thắng lợi chống ngoại xâm và giữ vững nền tự chủ quốc gia.

Thứ tư, trên bình diện giá trị nhân văn, sự dung hợp cởi mở giữa Nho -- Phật -- Đạo cùng tiếng nói phê phán lễ giáo phong kiến trong văn học chữ Nôm đã đánh dấu sự thức tỉnh của ý thức cá nhân. Con người hiện thực với nhu cầu hạnh phúc trần thế, đặc biệt là phẩm giá và quyền sống của người phụ nữ, lần đầu tiên trở thành mối quan tâm trung tâm của đời sống tinh thần, tạo bước chuyển quan trọng trong lịch sử văn hóa Việt Nam.

\subsubsection*{3.1.2. Những hạn chế mang tính lịch sử và rào cản thế giới quan giai cấp phong kiến}
\addcontentsline{toc}{subsubsection}{3.1.2. Những hạn chế mang tính lịch sử và rào cản thế giới quan giai cấp phong kiến}

Song song với những đóng góp tiến bộ, do bị quy định bởi tồn tại xã hội đương thời, tầng lớp nho sĩ thế kỷ XVIII vẫn mang những giới hạn lịch sử không thể vượt qua:

Trước hết là giới hạn về thế giới quan chính trị -- giai cấp. Mặc dù kinh tế hàng hóa thế kỷ XVIII đã có bước khởi sắc, xã hội Đại Việt về căn bản vẫn dựa trên phương thức sản xuất phong kiến tiểu nông phân tán, chưa xuất hiện phương thức sản xuất tư bản chủ nghĩa cũng như giai cấp tư sản dân tộc. Bởi vậy, tư duy chính trị của ngay cả những nho sĩ tiến bộ nhất vẫn nằm trong quỹ đạo quân chủ: lý tưởng cao nhất của họ chỉ dừng lại ở mô hình ``minh quân -- nhân chính'' (trông đợi một bậc vua sáng tôi hiền cứu đời), chứ chưa thể hình thành ý niệm về quyền làm chủ của nhân dân hay thiết chế nhà nước dân chủ pháp quyền.

Tiếp đến là giới hạn về trình độ nhận thức khoa học. Trong điều kiện nền sản xuất thủ công và sự vắng mặt của khoa học tự nhiên thực nghiệm hiện đại, những yếu tố duy vật và biện chứng trong học thuyết của các danh nho thế kỷ XVIII mới mang tính chất trực quan, tự phát và kinh nghiệm chủ nghĩa. Khi lý giải các hiện tượng xã hội phức tạp hoặc biến cố tự nhiên, tư duy của họ đôi lúc vẫn pha tạp thuyết thiên mệnh, sấm vĩ và ngũ hành cảm ứng.

Cuối cùng là sức ỳ bảo thủ của ý thức hệ phong kiến lạc hậu. Sự bám rễ sâu của luân lý tam cương và tâm lý đẳng cấp ``kẻ sĩ'' đã khiến một bộ phận không nhỏ trí thức quay lưng lại với phong trào quần chúng: kẻ cố chấp giữ tiết cực đoan đến mức tuẫn vong hoặc tiếp tay cho ngoại bang chống lại dân tộc; người lại thu mình trong lối sống ẩn dật tiêu cực, hoài cổ, gây lãng phí nguồn lực chất xám trước yêu cầu cải cách đất nước.

\subsection*{3.2. Ý nghĩa lịch sử và bài học kinh nghiệm trong việc phát huy vai trò đội ngũ trí thức giai đoạn hiện nay}
\addcontentsline{toc}{subsection}{3.2. Ý nghĩa lịch sử và bài học kinh nghiệm trong việc phát huy vai trò đội ngũ trí thức giai đoạn hiện nay}

\subsubsection*{3.2.1. Bài học về nguyên tắc thống nhất giữa lý luận với thực tiễn và việc rèn luyện bản lĩnh chính trị}
\addcontentsline{toc}{subsubsection}{3.2.1. Bài học về nguyên tắc thống nhất giữa lý luận với thực tiễn và việc rèn luyện bản lĩnh chính trị}

Từ sự thăng trầm của tầng lớp nho sĩ thế kỷ XVIII, bài học đầu tiên có giá trị thời sự trực tiếp là phải quán triệt nghiêm túc nguyên tắc thống nhất giữa lý luận với thực tiễn, kiên quyết loại bỏ bệnh hình thức và tệ sùng bái bằng cấp. Hệ lụy từ việc thương mại hóa tư cách trường thi thời Lê mạt cho thấy: một khi giáo dục tách rời thực tiễn và học vị biến thành công cụ mưu cầu danh lợi thay vì phản ánh năng lực thực chất, nền học thuật sẽ suy thoái từ gốc rễ.

Trong bối cảnh phát triển kinh tế tri thức và chuyển đổi số hiện nay, công tác giáo dục -- đào tạo và quản lý khoa học cần tiếp tục đấu tranh triệt để với bệnh thành tích, tình trạng ``học giả bằng thật'' và thiếu liêm chính học thuật. Kế thừa tinh thần thực học của tiền nhân, giá trị của người trí thức hiện đại phải được kiểm chứng thông qua năng lực sáng tạo công nghệ lõi, các công trình nghiên cứu có khả năng ứng dụng vào sản xuất, quản trị quốc gia và nâng cao chất lượng sống của nhân dân.

Gắn liền với năng lực chuyên môn là bài học về bản lĩnh chính trị và đạo đức trách nhiệm của người trí thức. Thực tiễn thế kỷ XVIII chứng minh rằng trí thức chỉ phát huy được tầm vóc lịch sử khi biết hòa mình vào dòng chảy tiến bộ của dân tộc; ngược lại, nếu đặt lợi ích cục bộ lên trên lợi ích chung hoặc giáo điều, xa rời nhân dân thì tất yếu sẽ rơi vào bế tắc. Trước những biến động đa chiều của tình hình thế giới và mặt trái cơ chế thị trường hiện nay, đội ngũ trí thức Việt Nam cần kiên định mục tiêu độc lập dân tộc gắn liền với chủ nghĩa xã hội trên nền tảng chủ nghĩa Mác -- Lênin, tư tưởng Hồ Chí Minh; nêu cao tinh thần phụng sự Tổ quốc, tích cực tham gia tư vấn, phản biện xã hội khách quan, khoa học nhằm hoàn thiện đường lối phát triển đất nước.

\subsubsection*{3.2.2. Bài học về quan điểm trọng dụng nhân tài và sự vận dụng trong đường lối của Đảng Cộng sản Việt Nam}
\addcontentsline{toc}{subsubsection}{3.2.2. Bài học về quan điểm trọng dụng nhân tài và sự vận dụng trong đường lối của Đảng Cộng sản Việt Nam}

Sự quy tụ đông đảo hiền tài Bắc Hà dưới ngọn cờ Tây Sơn để lại bài học quý giá về nghệ thuật chiêu hiền đãi sĩ trong quản trị quốc gia. Muốn thu hút và phát huy chất xám của đội ngũ trí thức, chủ thể lãnh đạo phải có thái độ cầu thị chân thành, bao dung, không định kiến hẹp hòi về xuất thân hay quá khứ, đồng thời mạnh dạn trao quyền, giao nhiệm vụ xứng tầm để người tài có không gian cống hiến thực chất.

Kế thừa truyền thống ``Hiền tài là nguyên khí của quốc gia'' trên lập trường mácxít, Đảng Cộng sản Việt Nam luôn xác định đội ngũ trí thức là lực lượng lao động sáng tạo đặc biệt quan trọng trong khối liên minh công nhân -- nông dân -- trí thức. Quan điểm nhất quán đó được cụ thể hóa thành tầm nhìn chiến lược tại Nghị quyết số 45-NQ/TW ngày 24/11/2023 của Hội nghị lần thứ tám Ban Chấp hành Trung ương Đảng Khóa XIII về ``Tiếp tục xây dựng và phát huy vai trò của đội ngũ trí thức đáp ứng yêu cầu phát triển đất nước nhanh và bền vững trong giai đoạn mới''.

Nghị quyết số 45-NQ/TW khẳng định xây dựng đội ngũ trí thức vững mạnh toàn diện chính là đầu tư trực tiếp cho bồi đắp nguyên khí quốc gia và phát triển bền vững. Để hiện thực hóa định hướng này, Đảng và Nhà nước tập trung hoàn thiện cơ chế tuyển dụng, trọng dụng, đãi ngộ đặc thù đối với các nhà khoa học đầu ngành và nhân tài trẻ dựa trên hiệu quả cống hiến; kiến tạo môi trường làm việc dân chủ, tôn trọng tự do học thuật gắn với đề cao đạo đức nghề nghiệp; đồng thời xây dựng hành lang pháp lý bảo vệ những cán bộ, trí thức dám nghĩ, dám làm, dám đổi mới sáng tạo vì lợi ích chung. Sự kết hợp hài hòa giữa chính sách đãi ngộ đúng đắn của thể chế với khát vọng dấn thân của bản thân người trí thức chính là động lực quyết định đưa đất nước vươn lên giàu mạnh, văn minh trong kỷ nguyên phát triển mới.

\newpage

% 7. Phần Kết luận
\section*{PHẦN KẾT LUẬN}
\addcontentsline{toc}{section}{PHẦN KẾT LUẬN}

Vận dụng thế giới quan và phương pháp luận của chủ nghĩa duy vật biện chứng và chủ nghĩa duy vật lịch sử để nghiên cứu đề tài ``Tìm hiểu về tầng lớp nho sĩ ở Việt Nam thế kỷ XVIII'', tiểu luận rút ra ba kết luận chủ yếu sau đây:

Thứ nhất, sự biến đổi của tầng lớp nho sĩ thế kỷ XVIII là minh chứng điển hình cho quy luật tồn tại xã hội quyết định ý thức xã hội và vị trí trung gian đặc thù của trí thức trong xã hội có giai cấp. Cuộc khủng hoảng toàn diện của thể chế Lê -- Trịnh, chúa Nguyễn, sự thương mại hóa khoa cử với hiện tượng ``Sinh đồ ba quan'' năm 1750, sự khởi sắc của kinh tế hàng hóa cùng làn sóng khởi nghĩa nông dân mà đỉnh cao là phong trào Tây Sơn đã làm sụp đổ thế độc tôn của Nho giáo chính thống. Chính sự biến động từ gốc rễ của cơ sở kinh tế -- xã hội đã buộc người nho sĩ phải xét lại các tín điều cũ để lựa chọn con đường hành xử mới.

Thứ hai, tầng lớp nho sĩ thế kỷ XVIII đã diễn ra sự phân hóa quyết liệt về lập trường chính trị, song hành cùng bước chuyển mình sâu rộng về học thuật và tư tưởng nhân văn. Trên bình diện chính trị, khác với bộ phận bảo thủ trói mình trong chữ ``trung quân'' giáo điều và bộ phận bất mãn lui về nông thôn mở trường tư thục giữ gìn văn hiến, những đại biểu ưu tú nhất như Ngô Thì Nhậm, Nguyễn Thiếp, Phan Huy Ích đã vượt qua rào cản đẳng cấp để dấn thân phò tá vua Quang Trung bảo vệ độc lập dân tộc; còn ở phương Nam, danh sư Võ Trường Toản với phương châm ``Tri ngôn dưỡng khí'' đã đào tạo nên lớp hiền tài làm rạng danh sĩ khí Gia Định. Trên bình diện học thuật, kẻ sĩ thế kỷ XVIII đã vượt thoát lối học từ chương Tống Nho để kiến tạo trào lưu thực học bách khoa mà đỉnh cao là tư tưởng ``Lý ở trong Khí'' và phương pháp khảo chứng thực địa của Lê Quý Đôn, kết hợp với xu hướng dung hợp Tam giáo và sự bùng nổ của chủ nghĩa nhân đạo trong văn học chữ Nôm.

Thứ ba, nhìn nhận khách quan theo nguyên tắc lịch sử -- cụ thể, dù chưa thể vượt qua giới hạn của phương thức sản xuất phong kiến tiểu nông và khuôn khổ lý tưởng ``minh quân -- nhân chính'', tầng lớp nho sĩ thế kỷ XVIII đã để lại di sản quý giá về tinh thần thực học, ý thức chủ quyền quốc gia và khát vọng giải phóng con người. Những kinh nghiệm lịch sử ấy mang ý nghĩa tham chiếu trực tiếp đối với hiện tại: kiên quyết bài trừ bệnh hư danh và bằng cấp hình thức, gắn kết chặt chẽ nghiên cứu lý luận với thực tiễn phát triển đất nước, bồi đắp bản lĩnh chính trị của người trí thức và thực thi hiệu quả chính sách trọng dụng hiền tài theo tinh thần Nghị quyết số 45-NQ/TW của Ban Chấp hành Trung ương Đảng Khóa XIII trong giai đoạn phát triển mới.

\newpage

% 8. Danh mục Tài liệu tham khảo
\section*{TÀI LIỆU THAM KHẢO}
\addcontentsline{toc}{section}{TÀI LIỆU THAM KHẢO}

\begin{enumerate}
    \item Đào Duy Anh (2014), \textit{Việt Nam văn hóa sử cương}, Quan Hải tùng thư xuất bản năm 1938; Nhà xuất bản Thế giới -- Công ty Văn hóa và Truyền thông Nhã Nam tái bản, Hà Nội.
    \item Ban Chấp hành Trung ương Đảng Cộng sản Việt Nam (2023), \textit{Nghị quyết số 45-NQ/TW ngày 24/11/2023 của Hội nghị lần thứ tám Ban Chấp hành Trung ương Đảng Khóa XIII về tiếp tục xây dựng và phát huy vai trò của đội ngũ trí thức đáp ứng yêu cầu phát triển đất nước nhanh và bền vững trong giai đoạn mới}, Hà Nội.
    \item Bộ Giáo dục và Đào tạo (2021), \textit{Giáo trình Triết học Mác -- Lênin} (Dành cho bậc đào tạo đại học và sau đại học hệ không chuyên lý luận chính trị), Nhà xuất bản Chính trị quốc gia Sự thật, Hà Nội.
    \item Lê Quý Đôn (1977), \textit{Vân đài loại ngữ} (Bản dịch của Viện Sử học), Nhà xuất bản Văn hóa Thông tin, Hà Nội.
    \item Lê Quý Đôn (1977), \textit{Phủ biên tạp lục} (Bản dịch của Viện Sử học), Nhà xuất bản Khoa học Xã hội, Hà Nội.
    \item Trần Văn Giàu (1973), \textit{Sự phát triển của tư tưởng ở Việt Nam từ thế kỷ XIX đến Cách mạng Tháng Tám} (Tập 1: Hệ ý thức phong kiến và sự thất bại của nó trước các nhiệm vụ lịch sử), Nhà xuất bản Khoa học Xã hội, Hà Nội.
    \item Phan Huy Lê, Chu Thiên, Vương Hoàng Tuyên, Đinh Xuân Lâm (1965), \textit{Lịch sử chế độ phong kiến Việt Nam} (Tập III -- Thời kỳ khủng hoảng và phong trào nông dân khởi nghĩa thế kỷ XVIII), Nhà xuất bản Giáo dục, Hà Nội.
    \item Ngô Thì Nhậm (2001), \textit{Trúc Lâm tông chỉ nguyên thanh} (Bản dịch của Viện Nghiên cứu Hán Nôm), Nhà xuất bản Hà Nội, Hà Nội.
    \item Nguyễn Khắc Thuần (1998), \textit{Đại cương lịch sử văn hóa Việt Nam} (Tập 4), Nhà xuất bản Giáo dục, Hà Nội.
    \item Nguyễn Tài Thư (Chủ biên, 1993), \textit{Lịch sử tư tưởng Việt Nam} (Tập 1), Viện Triết học -- Viện Khoa học Xã hội Việt Nam, Nhà xuất bản Khoa học Xã hội, Hà Nội.
    \item Đặng Việt Bích (2015), ``Quang Trung trong tâm trí kẻ sĩ Bắc Hà'', \textit{Tạp chí Văn học Nghệ thuật}.
    \item Quách Hiền (2018), ``Tư học Việt Nam thế kỷ XVIII'' (Khảo sát quá trình chuyển dịch của các danh nho về làng xã mở trường tư thục), chuyên trang Nghiên Cứu Lịch Sử.
    \item Nguyễn Thị Hồng Phương (2026), ``Nhân sinh quan trong tư tưởng triết học của Lê Quý Đôn -- Ý nghĩa đối với xây dựng con người Việt Nam hiện nay'', \textit{Tạp chí Nghiên cứu Triết học}, Số 31 (2-2026), tr. 62--69.
    \item Nguyễn Kim Sơn (2004), ``Những nhân tố làm khởi phát khuynh hướng khảo chứng học thế kỷ XVIII -- nửa đầu thế kỷ XIX'', \textit{Kỷ yếu Thông báo Hán Nôm học}, Viện Nghiên cứu Hán Nôm -- Viện Hàn lâm Khoa học xã hội Việt Nam.
    \item Cao Văn Thức (2016), ``Tản mạn về học vị Sinh Đồ và chuyện Sinh Đồ ba quan'', chuyên trang Nghiên Cứu Lịch Sử.
    \item Sở Giáo dục và Đào tạo Thành phố Hồ Chí Minh và Thư viện tỉnh Đồng Nai (2020), Chuyên đề: ``Danh sư Võ Trường Toản -- Bậc Sùng Đức xử sĩ và hào khí Đồng Nai xứ Gia Định thế kỷ XVIII''.
\end{enumerate}

\newpage

\end{document}

